%% =====================================================================
%% Canovaccio della discussione NDS (prof. Baggio), 26/11/2026, ore 9:30
%% Solo il discorso: per ogni slide le frasi da dire, nell'ordine della slide.
%% =====================================================================
%% COME SI AGGIORNA
%% 1. Ogni slide e' un blocco
%%       \begin{slide}{<label del frame>}{<titolo>}{<tempo>}
%%         \begin{discorso}
%%           \item prima frase
%%           \item seconda frase
%%         \end{discorso}
%%       \end{slide}
%%    La <label> e' quella del frame in main.tex, es. [label=m-bridge].
%%    Frame a due passi: \begin{slidedue}{...}{...}{...}; \clic segna il passo 2.
%%    Pagine senza label (titolo, outline, Thank you): \begin{slidepag}{<pagina>}{...}{...}.
%% 2. Miniature, numeri di slide e di pagina NON si scrivono a mano: vengono
%%    da ../main.pdf e ../main.aux. Se cambi la presentazione, ricompila prima
%%    main.tex e poi questo file (pdflatex canovaccio.tex, due volte).
%% 3. Nel testo: \fr{label} = numero della slide; \pg{label} = pagina del PDF.
%% 4. Regole del discorso: niente numeri di teoremi, corollari o casi (si dice il
%%    contenuto o il nome); si dice solo cio' che e' sulla slide; termini tecnici in
%%    inglese come sulle slide (consensus, disagreement, prejudice, susceptibility,
%%    attachment, augmented, social power, influence centrality, ...).
%% 5. Domande probabili, termini da evitare e cose nostre: file domande.tex.
%% =====================================================================
\documentclass[a4paper,10pt]{article}
\usepackage[T1]{fontenc}
% babel in italiano: con TeX Live completo usa italian.ldf, altrimenti il file .ini di babel
\IfFileExists{italian.ldf}{\usepackage[italian]{babel}}%
  {\usepackage{babel}\babelprovide[import,main]{italian}}
\usepackage{lmodern}
\usepackage{microtype}
\usepackage[a4paper,margin=1.5cm,top=1.4cm,bottom=1.7cm]{geometry}
\usepackage{amsmath,amssymb}
\usepackage{graphicx}
\usepackage{xcolor}
\usepackage{enumitem}
\usepackage{refcount}
\usepackage{xr-hyper}
\usepackage[hidelinks]{hyperref}
\externaldocument[P-]{../main}  % la presentazione nuova
\externaldocument[O-]{../../presentation_v8/main}  % la presentazione precedente
\pdfsuppresswarningpagegroup=1

% ---------------------------------------------------------------------
% Riferimenti alla presentazione
\newcommand{\presentazione}{../main.pdf}
\newcommand{\pg}[1]{\getpagerefnumber{P-#1}}  % pagina del PDF
\newcommand{\fr}[1]{\getrefnumber{P-#1}}      % numero della slide
\newcommand{\uno}{\mathbf{1}}

% ---------------------------------------------------------------------
% Impaginazione
\setlength{\parindent}{0pt}
\setlength{\parskip}{2pt}
\raggedbottom
\definecolor{unipd}{RGB}{155,0,20}

\newcommand{\mini}[1]{{\setlength{\fboxsep}{0pt}%
  \fbox{\includegraphics[page=#1,width=\dimexpr\linewidth-2\fboxrule\relax]{\presentazione}}}}
\newcommand{\testata}[4]{% numero, pagine, titolo, tempo
  {\color{unipd}\bfseries Slide #1 \textperiodcentered\ #3}\hfill{\small\itshape #4}\par
  {\scriptsize\color{gray} pagina #2 del PDF}\par\smallskip}

% Blocco di una slide: miniatura a sinistra, discorso a destra
\newenvironment{slide}[3]{%
  \par\addvspace{12pt}\noindent
  \begin{minipage}[t]{0.36\textwidth}\vspace{0pt}\mini{\pg{#1}}\end{minipage}\hfill
  \begin{minipage}[t]{0.61\textwidth}\vspace{0pt}%
  \testata{\fr{#1}}{\pg{#1}}{#2}{#3}}%
 {\end{minipage}\par}
% Frame a due passi: le due miniature una sotto l'altra
\newenvironment{slidedue}[3]{%
  \par\addvspace{12pt}\noindent
  \begin{minipage}[t]{0.36\textwidth}\vspace{0pt}\mini{\pg{#1<1>}}\par\smallskip
    \mini{\pg{#1<2>}}\end{minipage}\hfill
  \begin{minipage}[t]{0.61\textwidth}\vspace{0pt}%
  \testata{\fr{#1}}{\pg{#1<1>}--\pg{#1<2>}}{#2}{#3}}%
 {\end{minipage}\par}
% ADDED (2026-11-18): slide con tre passi
\newenvironment{slidetre}[3]{%
  \par\addvspace{12pt}\noindent
  \begin{minipage}[t]{0.36\textwidth}\vspace{0pt}\mini{\pg{#1<1>}}\par\smallskip
    \mini{\pg{#1<2>}}\par\smallskip\mini{\pg{#1<3>}}\end{minipage}\hfill
  \begin{minipage}[t]{0.61\textwidth}\vspace{0pt}%
  \testata{\fr{#1}}{\pg{#1<1>}--\pg{#1<3>}}{#2}{#3}}%
 {\end{minipage}\par}
% Pagine senza label: si indica la pagina del PDF
\newenvironment{slidepag}[3]{%
  \par\addvspace{12pt}\noindent
  \begin{minipage}[t]{0.36\textwidth}\vspace{0pt}\mini{#1}\end{minipage}\hfill
  \begin{minipage}[t]{0.61\textwidth}\vspace{0pt}%
  {\color{unipd}\bfseries #2}\hfill{\small\itshape #3}\par
  {\scriptsize\color{gray} pagina #1 del PDF}\par\smallskip}%
 {\end{minipage}\par}

% Il discorso: una frase (o due) per punto
\newenvironment{discorso}{\begin{itemize}[nosep,leftmargin=1em,label={\color{gray}\textbullet},itemsep=2pt]}{\end{itemize}}
\newcommand{\clic}{{\color{unipd}\scriptsize[passo 2]}\ }
\newcommand{\parte}[2]{\section*{#1\hfill{\normalsize\itshape #2}}}


% testo tolto: barrato parola per parola (senza pacchetti esterni)
\newsavebox{\barratobox}
\newcommand{\barrato}[1]{\leavevmode\sbox{\barratobox}{#1}%
  \rlap{\rule[0.55ex]{\wd\barratobox}{0.5pt}}\usebox{\barratobox}}
\newcommand{\vecchia}{../../presentation_v8/main.pdf}
\newcommand{\pgO}[1]{\getpagerefnumber{O-#1}}
\newcommand{\frO}[1]{\getrefnumber{O-#1}}
\newcommand{\miniO}[1]{{\setlength{\fboxsep}{0pt}%
  \fbox{\includegraphics[page=#1,width=\dimexpr\linewidth-2\fboxrule\relax]{\vecchia}}}}
\newcommand{\etichetta}[2]{\,{\setlength{\fboxsep}{1.5pt}\colorbox{#1}{\scriptsize\bfseries\color{white}#2}}}
\newcommand{\nuova}{\etichetta{blue}{NUOVA}}
\newcommand{\spostata}{\etichetta{orange}{SPOSTATA}}
\newcommand{\divisa}{\etichetta{teal}{DIVISA}}
\newcommand{\toltatesto}[1]{{\color{red}\barrato{#1}}}
\begin{document}
\begin{center}
  {\LARGE\bfseries Preview delle modifiche}\\[3pt]
  {\large Presentazione e canovaccio, discussione NDS del 26/11}\\[2pt]
  {\small\color{gray} Versione del \today. Confronto con la versione precedente: presentazione e canovaccio applicati il 18/11 (v8).}
\end{center}

\textbf{Legenda.} {\color{blue}Blu}: testo nuovo. {\color{red}\barrato{Rosso barrato}}: testo
tolto (le formule tolte sono solo in rosso). \nuova: slide o blocco nuovo.
\spostata: slide spostata. \divisa: slide divisa in due.
\section*{Slide modificate: prima e dopo}
Solo le slide cambiate, nel nuovo ordine. In tutto 33 slide (prima 31): sono nuove «All agents P-dependent» e «Example networks: bidirectional ring». Le note del 19/11 sono tutte qui, tranne le slide sulla social power, che non ho toccato (vedi in fondo).
\par\addvspace{14pt}\noindent\begin{minipage}{\textwidth}
{\color{unipd}\bfseries Slide \fr{m-fdg-rosetta} \textperiodcentered\ Which outcome: the Rosetta stone}\spostata\par\smallskip
{\small Spostata dopo «Disagreement and stubbornness», come nel report (§2.1.4). Ora l'ordine è: The FDG model, Consensus and social power, Disagreement and stubbornness, Rosetta stone, Two networks. Tolta la riga «Focus of this work…»; al suo posto la frase del report: la traduzione fra $A$ e $G$ viene dal Theorem 1.11 nei casi (i)--(ii) e dai Theorems 2.2 e 2.4 nei casi (iii)--(iv). Corollary 1.12 più stretto, per non finire sotto i pulsanti. Nuovo pulsante «Theorem 1.11», che porta al backup con il teorema; da lì si torna alla Rosetta stone. La frase sul focus va nel discorso di «Consensus and social power»: su quella slide non c'è spazio.}\par\smallskip
\noindent\begin{minipage}[t]{0.36\textwidth}\vspace{0pt}{\scriptsize\color{gray} Prima}\par
\miniO{\pgO{m-fdg-rosetta}}\par\smallskip
\end{minipage}\hfill\begin{minipage}[t]{0.61\textwidth}\vspace{0pt}{\scriptsize\color{gray} Dopo}\par
\mini{\pg{m-fdg-rosetta}}\par\smallskip
\end{minipage}\par
\end{minipage}\par
\par\addvspace{14pt}\noindent\begin{minipage}{\textwidth}
{\color{unipd}\bfseries Slide \fr{m-networks} \textperiodcentered\ Two networks, one agent}\par\smallskip
{\small Due passi. Passo 1: le reti; la parametrizzazione con la frase concordata («Agent 1 gives weight $t$ to itself and scales its original weights by $1-t$ … with $a_1^\top$ the row of agent 1 at $t = 0$. The other rows do not change»); la riga dell'agente 1 scritta per le due reti; $x(0)$. Passo 2: il caso $t = 0$ e la tabella. Didascalia: «Unlabelled edges: weight 1 in (a), 1/2 in (b)»: sull'anello 2 e 6 ascoltano l'agente 1 con peso 1/2 per ogni $t$. Il self-loop tratteggiato dell'agente 1 è più largo: prima si vedeva appena. Colonna destra in \texttt{footnotesize} e un po' più larga, per farci stare tutto.}\par\smallskip
\noindent\begin{minipage}[t]{0.36\textwidth}\vspace{0pt}{\scriptsize\color{gray} Prima}\par
\miniO{\pgO{m-networks}}\par\smallskip
\end{minipage}\hfill\begin{minipage}[t]{0.61\textwidth}\vspace{0pt}{\scriptsize\color{gray} Dopo}\par
\mini{\pg{m-networks<1>}}\par\smallskip
\mini{\pg{m-networks<2>}}\par\smallskip
\end{minipage}\par
\end{minipage}\par
\par\addvspace{14pt}\noindent\begin{minipage}{\textwidth}
{\color{unipd}\bfseries Slide \fr{m-fj-conv}--\fr{m-fj-conv2} \textperiodcentered\ Convergence in the FJ model; Convergence with P-independent agents}\par\smallskip
{\small Ordine del report (§3.1.2): $V$ non compare più prima di essere definita. Slide \fr{m-fj-conv}: definizioni, Definition 3.1, poi la decomposizione (6)--(7) con $W^{21} = 0$, $\Lambda^{22} = I$ e il significato di $W^{22}$: è row-stochastic, quindi (7) è un modello FDG fra i P-independent. Slide \fr{m-fj-conv2}: solo il Theorem 3.2, che definisce $V$ in generale. Nel riquadro una frase in più, dal report: (7) converge se e solo se tutti i sink del condensation digraph dei P-independent sono aperiodici. Così si vede dove serve l'aperiodicità.}\par\smallskip
\noindent\begin{minipage}[t]{0.36\textwidth}\vspace{0pt}{\scriptsize\color{gray} Prima}\par
\miniO{\pgO{m-fj-conv}}\par\smallskip
\miniO{\pgO{m-fj-conv2}}\par\smallskip
\end{minipage}\hfill\begin{minipage}[t]{0.61\textwidth}\vspace{0pt}{\scriptsize\color{gray} Dopo}\par
\mini{\pg{m-fj-conv}}\par\smallskip
\mini{\pg{m-fj-conv2}}\par\smallskip
\end{minipage}\par
\end{minipage}\par
\par\addvspace{14pt}\noindent\begin{minipage}{\textwidth}
{\color{unipd}\bfseries Slide \fr{m-fj-conv3} \textperiodcentered\ All agents P-dependent}\nuova\par\smallskip
{\small Il Corollary 3.3 come caso $r = n$ del Theorem 3.2, poi $x(\infty) = Vu$ (9) e l'intuizione del report: $Vu$ è un equilibrio, $e(k+1) = \Lambda W e(k)$ (10), «No aperiodicity is needed». Il pulsante «Stability» è qui, e il backup torna qui.}\par\smallskip
\noindent\begin{minipage}[t]{0.36\textwidth}\vspace{0pt}{\scriptsize\color{gray} Prima}\par
{\small ---}
\end{minipage}\hfill\begin{minipage}[t]{0.61\textwidth}\vspace{0pt}{\scriptsize\color{gray} Dopo}\par
\mini{\pg{m-fj-conv3}}\par\smallskip
\end{minipage}\par
\end{minipage}\par
\par\addvspace{14pt}\noindent\begin{minipage}{\textwidth}
{\color{unipd}\bfseries Slide \fr{m-fj-final} \textperiodcentered\ Final opinions and the augmented graph}\par\smallskip
{\small Tolta la frase sul disagreement (caso (iv), Theorem 2.4, limite (9)): la figura ha un solo prejudice, e il teorema del disagreement passa nel discorso di «Back to the question». Resta il fatto che serve subito dopo, per la influence centrality (il modello augmented ha almeno $n$ sink): «Each virtual agent is stubborn: a sink of the augmented graph, aperiodic because of its self-loop». Che i grafi sono semplificati (12 nodi, 7 disegnati) è detto solo nel canovaccio.}\par\smallskip
\noindent\begin{minipage}[t]{0.36\textwidth}\vspace{0pt}{\scriptsize\color{gray} Prima}\par
\miniO{\pgO{m-fj-final}}\par\smallskip
\end{minipage}\hfill\begin{minipage}[t]{0.61\textwidth}\vspace{0pt}{\scriptsize\color{gray} Dopo}\par
\mini{\pg{m-fj-final}}\par\smallskip
\end{minipage}\par
\end{minipage}\par
\par\addvspace{14pt}\noindent\begin{minipage}{\textwidth}
{\color{unipd}\bfseries Slide \fr{m-fj-graphs} \textperiodcentered\ Example networks: directed cycle}\par\smallskip
{\small Il titolo dice la rete e sta su una riga (prima «The networks of the examples»). Tolto «the bidirectional ring in the same way» dalla didascalia. Self-loop tratteggiato di (a) più largo.}\par\smallskip
\noindent\begin{minipage}[t]{0.36\textwidth}\vspace{0pt}{\scriptsize\color{gray} Prima}\par
\miniO{\pgO{m-fj-graphs}}\par\smallskip
\end{minipage}\hfill\begin{minipage}[t]{0.61\textwidth}\vspace{0pt}{\scriptsize\color{gray} Dopo}\par
\mini{\pg{m-fj-graphs}}\par\smallskip
\end{minipage}\par
\end{minipage}\par
\par\addvspace{14pt}\noindent\begin{minipage}{\textwidth}
{\color{unipd}\bfseries Slide \fr{m-fj-graphs-ring} \textperiodcentered\ Example networks: bidirectional ring}\nuova\par\smallskip
{\small Le stesse quattro reti sull'anello. (a)--(b): l'agente 1 ascolta 2 e 6 con $\frac{1-t}{2}$ ciascuno, con self-loop $t$ in (a) e ingresso $t$ da $u_1$ in (b); tutti gli altri archi pesano $1/2$, anche quelli con cui 2 e 6 ascoltano l'agente 1. (c)--(d): self-loop o prejudice $0.2$, ogni arco $0.4$.}\par\smallskip
\noindent\begin{minipage}[t]{0.36\textwidth}\vspace{0pt}{\scriptsize\color{gray} Prima}\par
{\small ---}
\end{minipage}\hfill\begin{minipage}[t]{0.61\textwidth}\vspace{0pt}{\scriptsize\color{gray} Dopo}\par
\mini{\pg{m-fj-graphs-ring}}\par\smallskip
\end{minipage}\par
\end{minipage}\par
\par\addvspace{14pt}\noindent\begin{minipage}{\textwidth}
{\color{unipd}\bfseries Slide \fr{m-fj-networks} \textperiodcentered\ FDG and FJ on the same networks}\par\smallskip
{\small Seconda riga scritta per le due reti. Ciclo: $x_1(k+1) = t\,x_1(k) + (1-t)\,x_2(k)$; anello: $t\,x_1(k) + (1-t)\,(x_2(k)+x_6(k))/2$. Nel FJ lo stesso, con $t\,u_1$ al posto di $t\,x_1(k)$. Sotto la tabella: nella terza riga ogni arco pesa 0.8 sul ciclo e 0.4 sull'anello.}\par\smallskip
\noindent\begin{minipage}[t]{0.36\textwidth}\vspace{0pt}{\scriptsize\color{gray} Prima}\par
\miniO{\pgO{m-fj-networks}}\par\smallskip
\end{minipage}\hfill\begin{minipage}[t]{0.61\textwidth}\vspace{0pt}{\scriptsize\color{gray} Dopo}\par
\mini{\pg{m-fj-networks}}\par\smallskip
\end{minipage}\par
\end{minipage}\par
\par\addvspace{14pt}\noindent\begin{minipage}{\textwidth}
{\color{unipd}\bfseries Slide \fr{m-state-input}, passo 2 \textperiodcentered\ Stubborn agent: state or input}\par\smallskip
{\small Qui il significato di $x(\infty) = Vu$, sulla perturbazione. Nel FDG il limite $x_1(0)\mathbf{1}_6$ diventa $0.5\,\mathbf{1}_6$; nel FJ $Vu$ non contiene $x(0)$, quindi il limite non cambia. Nel canovaccio c'è anche la ricorsione $\delta(k+1) = \Lambda W\,\delta(k)$ del report.}\par\smallskip
\noindent\begin{minipage}[t]{0.36\textwidth}\vspace{0pt}{\scriptsize\color{gray} Prima}\par
\miniO{\pgO{m-state-input<2>}}\par\smallskip
\end{minipage}\hfill\begin{minipage}[t]{0.61\textwidth}\vspace{0pt}{\scriptsize\color{gray} Dopo}\par
\mini{\pg{m-state-input<2>}}\par\smallskip
\end{minipage}\par
\end{minipage}\par
\par\smallskip\noindent \textbf{Invariate.} «Consensus and social power»: la frase sul focus è solo nel discorso. Le slide sulla social power («Social power grows gradually», «Why the same social power»): non le ho toccate. Dimmi cosa non ti convince e le rifaccio.\par
\par\smallskip\noindent \textbf{Backup.} «Asymptotic stability» torna alla slide \fr{m-fj-conv3}. «Primitive, substochastic» ha anche il pulsante «Back to the Rosetta stone».\par
\par\smallskip\noindent \textbf{Una nota sul report.} La frase sulla Rosetta stone attribuisce al Theorem 1.11 anche la caratterizzazione delle matrici irriducibili. Nel report quella è l'enunciato prima del teorema: il Theorem 1.11 riguarda solo le primitive. Sulla slide ho lasciato la frase come nel report; la risposta pronta è in domande.tex.\par
\par\smallskip\noindent \textbf{domande.tex.} Aggiornati i numeri di slide: Stability e intuizione su \fr{m-fj-conv3}; la lettura del modello augmented come caso (iv) su \fr{m-back-question} e \fr{m-six-sinks}. Aggiunta la domanda sul Theorem 1.11.\par
\clearpage
\section*{Canovaccio: modifiche rispetto alla versione precedente}
Cambiano: le transizioni alla fine di «The FDG model» e di «Disagreement and stubbornness»; la Rosetta stone dopo il disagreement, come riassunto di tutti i casi; il focus sul consensus nel discorso di «Consensus and social power»; «Two networks, one agent» in due passi; la convergenza FJ in tre blocchi; «Final opinions» senza il teorema del disagreement, che passa a «Back to the question», e con i grafi semplificati detti a voce; il blocco nuovo dell'anello; la tabella per le due reti; il significato di $Vu$ in «State or input». Tempi per parte aggiornati: in tutto 29:05.
\parte{Parte 1 — Introduzione}{3:00}
\begin{slidepag}{1}{Titolo}{0:15}
  \begin{discorso}
    \item Buongiorno. Presento il mio progetto su stubbornness e social power nella dinamica delle opinioni.
    \item Studio due modelli, il modello di French–Harary–DeGroot e il modello di Friedkin–Johnsen: come cambia il loro comportamento quando gli agenti resistono agli altri, e quanto pesa ciascun agente nelle opinioni finali.
  \end{discorso}
\end{slidepag}

\begin{slidepag}{2}{Slide 2 \textperiodcentered\ Outline}{0:15}
  \begin{discorso}
    \item La presentazione ha quattro parti: introduzione, modello FDG, modello FJ, conclusioni.
    \item Alcune slide hanno dei pulsanti in basso a destra: aprono le slide di backup, con definizioni e derivazioni. Le uso solo se vuole entrare nel dettaglio.
  \end{discorso}
\end{slidepag}

\begin{slide}{m-question}{Stubbornness}{1:00}
  \begin{discorso}
    \item Una rete sociale è un digrafo pesato: gli agenti sono i nodi, gli archi portano l'influenza, e ogni agente aggiorna la sua opinione facendo la media delle opinioni che ascolta.
    \item Un agente è stubborn se la sua opinione resta invariata indipendentemente dalle opinioni degli altri: è la definizione del tutorial di Proskurnikov e Tempo.
    \item È il caso estremo di due modi graduati di resistere agli altri: dare peso alla propria opinione attuale, oppure restare attaccati a un'opinione fissa propria, un prejudice.
    \item Il primo è il modo del modello FDG, il secondo del modello FJ: li introduco tra poco.
    \item A entrambi i modelli, sulle stesse reti, faccio tre domande. Le opinioni convergono? Convergono alla stessa opinione per tutti, il consensus, o a opinioni diverse, il disagreement? E quanto pesa l'opinione iniziale di ciascun agente in quelle finali?
    \item Prima di tutto, da dove vengono i due modelli.
  \end{discorso}
\end{slide}

\begin{slide}{m-models}{Models of opinion dynamics}{0:40}
  \begin{discorso}
    \item Il quadro storico è quello del tutorial di Proskurnikov e Tempo, la fonte principale del lavoro, insieme alle slide del corso e al libro di Bullo.
    \item A sinistra i modelli che fanno solo medie: in tempo discreto il modello di French, Harary e DeGroot; in tempo continuo quello di Abelson.
    \item Per i gruppi reali il consensus è una previsione troppo forte: spesso restano divisi.
    \item Per questo il modello di Friedkin e Johnsen, negli anni Novanta, aggiunge alle medie opinioni esterne costanti; è stato validato sperimentalmente su gruppi piccoli e medi. In tempo continuo il corrispondente è il modello di Taylor.
    \item In questo lavoro tratto il modello FDG, visto nel corso, e il modello FJ, che è nuovo; Abelson e Taylor restano solo nella panoramica.
    \item Richiamo solo le nozioni del corso che mi servono.
  \end{discorso}
\end{slide}

\begin{slide}{m-prelim}{Preliminaries}{0:50}
  \begin{discorso}
    \item Sono nozioni del corso, quindi solo i punti che userò.
    \item Un arco da $i$ a $j$, con peso $a_{ij}$, vuol dire che l'agente $i$ è influenzato dall'agente $j$: è la convenzione del corso; il tutorial disegna gli archi al contrario.
    \item Un digrafo fortemente connesso è periodico se un intero maggiore di uno divide la lunghezza di ogni ciclo semplice; un self-loop lo rende aperiodico.
    \item Nel condensation digraph ogni componente fortemente connessa diventa un nodo; un sink è un gruppo influenzato solo da sé stesso. Saranno i sink a decidere l'esito.
    \item Per le matrici: semi-convergente vuol dire che il limite delle potenze esiste; convergente, o Schur stabile, che il limite è zero, cioè raggio spettrale minore di uno.
    \item Infine: irriducibile equivale a fortemente connesso, primitiva a fortemente connesso e aperiodico.
    \item Ora il primo modello.
  \end{discorso}
\end{slide}

\parte{Parte 2 — Il modello FDG}{8:40 {\color{red}\barrato{8:10}}}
\begin{slide}{m-fdg-model}{The FDG model}{0:50}
  \begin{discorso}
    \item Nel modello FDG ogni agente sostituisce la sua opinione con una media pesata delle opinioni che ascolta.
    \item In forma vettoriale $x(k+1)=Ax(k)$, con $A$ la matrice di adiacenza pesata, row-stochastic: è un sistema di media a tempo discreto.
    \item Iterando, $x(k)=A^kx(0)$: il comportamento nel lungo periodo dipende dalle potenze di $A$. Una matrice row-stochastic ha sempre l'autovalore 1, e raggio spettrale 1.
    \item A destra la definizione: il modello è convergente se, per ogni condizione iniziale, il limite esiste; raggiunge un consensus se, per ogni condizione iniziale, le opinioni finali sono tutte uguali.
    \item Gli esiti possibili sono quindi tre: nessuna convergenza, cioè $A$ non semi-convergente; consensus; e disagreement, cioè il modello converge ma per qualche condizione iniziale le opinioni finali non sono tutte uguali.
    \item {\color{blue}Quale dei tre si ha lo dicono due teoremi del corso, uno per il consensus e uno per il disagreement. Comincio dal consensus.}
    \item {\color{red}\barrato{Quale} \barrato{dei} \barrato{tre} \barrato{si} \barrato{ha} \barrato{si} \barrato{legge} \barrato{in} \barrato{tre} \barrato{modi} \barrato{equivalenti.}}
  \end{discorso}
\end{slide}

\begin{slide}{m-fdg-consensus}{Consensus and social power}{0:55 {\color{red}\barrato{0:50}}}
  \begin{discorso}
    \item Il teorema del consensus del corso dà tre condizioni equivalenti.
    \item Sullo spettro: l'autovalore 1 è semplice e gli altri hanno modulo minore di uno. Sulla matrice: le potenze di $A$ convergono a $\uno w^\top$. Sul grafo: esiste un nodo globally reachable, e il sottografo dei nodi globally reachable è aperiodico.
    \item In questo caso $w$ è l'autovettore sinistro dominante di $A$, con $w_i$ positivo se e solo se il nodo $i$ è globally reachable, e il valore di consensus è $w^\top x(0)$.
    \item A destra: $w_i$ è la social power dell'agente $i$, il peso della sua opinione iniziale nel valore di consensus.
    \item Richiede il consensus, perché solo allora le righe del limite di $A^k$ sono tutte uguali a $w^\top$.
    \item {\color{blue}Per questo il lavoro si concentra sul consensus: è il caso in cui il peso di ciascun agente si può misurare.}
    \item Sotto, l'average consensus: se $A$ è doubly-stochastic e vale la condizione sul grafo, $w$ è uniforme e il consensus è la media delle opinioni iniziali.
    \item Nei termini del tutorial: il modello converge se e solo se $A$ è regular, e raggiunge il consensus se e solo se $A$ è fully regular.
    \item E se il consensus non c'è?
  \end{discorso}
\end{slide}

\begin{slide}{m-fdg-disagreement}{Disagreement and stubbornness}{0:40}
  \begin{discorso}
    \item Se il condensation digraph ha almeno due sink, il teorema del disagreement dà tre condizioni equivalenti: l'autovalore 1 è semi-semplice, con molteplicità pari al numero di sink; $A$ è semi-convergente; ogni sink è aperiodico.
    \item Allora ogni sink raggiunge un consensus al suo interno, e ogni agente fuori dai sink converge a una combinazione convessa dei valori dei sink.
    \item Le righe del limite di $A^k$ sono diverse: nessun vettore fa la parte di $w$, e la social power non è definita.
    \item A destra la definizione di stubborn del tutorial. Nel modello FDG l'agente $i$ è stubborn se e solo se $a_{ii}=1$: è un sink fatto di un solo nodo, aperiodico grazie al self-loop.
    \item Se tutti gli agenti sono influenzati da lui, direttamente o attraverso altri, è l'unico nodo globally reachable, e la sua social power è 1.
    \item {\color{blue}Riassumo tutti i casi nella Rosetta stone.}
    \item {\color{red}\barrato{Ora} \barrato{le} \barrato{due} \barrato{reti} \barrato{su} \barrato{cui} \barrato{confronto} \barrato{i} \barrato{due} \barrato{modelli.}}
  \end{discorso}
\end{slide}

\begin{slide}{m-fdg-rosetta}{Which outcome: the Rosetta stone\spostata}{0:45 {\color{red}\barrato{0:40}}}
  \begin{discorso}
    \item L'esito si legge sullo spettro di $A$, sulla matrice o sul digrafo: la Rosetta stone del corso collega i tre punti di {\color{red}\barrato{vista,}} {\color{blue}vista} e {\color{red}\barrato{ogni} \barrato{riga} \barrato{della} \barrato{tabella} \barrato{è} \barrato{un} \barrato{caso.}} {\color{blue}riassume tutti i casi, uno per riga.}
    \item {\color{blue}Nelle prime due righe il digrafo è fortemente connesso. Se è anche aperiodico, $A$ è primitiva e si ha consensus; se è periodico, nessuna convergenza. Qui la traduzione fra matrice e digrafo viene dalla caratterizzazione delle matrici irriducibili e primitive.}
    \item {\color{blue}Nella terza e nella quarta riga il digrafo non è fortemente connesso, e la traduzione viene dai due teoremi appena visti. Un solo sink, aperiodico: consensus. Almeno due sink, tutti aperiodici: disagreement.}
    \item {\color{red}\barrato{Sul} \barrato{digrafo} \barrato{conta} \barrato{solo} \barrato{il} \barrato{condensation} \barrato{digraph:} \barrato{quanti} \barrato{sink} \barrato{ha,} \barrato{e} \barrato{se} \barrato{sono} \barrato{periodici.}}
    \item {\color{red}\barrato{Se} \barrato{c'è} \barrato{un} \barrato{solo} \barrato{sink,} \barrato{ed} \barrato{è} \barrato{aperiodico,} \barrato{si} \barrato{ha} \barrato{consensus:} \barrato{è} \barrato{il} \barrato{caso} \barrato{di} \barrato{una} \barrato{rete} \barrato{fortemente} \barrato{connessa} \barrato{e} \barrato{aperiodica,} \barrato{o} \barrato{di} \barrato{un} \barrato{unico} \barrato{sink} \barrato{aperiodico.}}
    \item {\color{red}\barrato{Se} \barrato{ci} \barrato{sono} \barrato{almeno} \barrato{due} \barrato{sink,} \barrato{tutti} \barrato{aperiodici,} \barrato{si} \barrato{ha} \barrato{disagreement.} \barrato{Se}} {\color{blue}L'ultima riga, lasciata vuota nelle slide del corso, viene dal tutorial: se} almeno un sink è periodico, il modello non converge.
    \item In basso il fatto che mi servirà di più: se il digrafo {\color{blue}ha almeno due nodi,} è fortemente connesso e ha almeno un self-loop, $A$ è primitiva.
    \item {\color{blue}Ora le due reti su cui confronto i due modelli.}
    \item {\color{red}\barrato{Il} \barrato{lavoro} \barrato{si} \barrato{concentra} \barrato{sul} \barrato{consensus,} \barrato{perché} \barrato{è} \barrato{il} \barrato{caso} \barrato{in} \barrato{cui} \barrato{il} \barrato{peso} \barrato{di} \barrato{ciascun} \barrato{agente} \barrato{si} \barrato{può} \barrato{misurare.} \barrato{Al} \barrato{disagreement} \barrato{dedico} \barrato{la} \barrato{slide} \barrato{dopo.}}
  \end{discorso}
\end{slide}

\begin{slidedue}{m-networks}{Two networks, one agent}{1:15 {\color{red}\barrato{0:55}}}
  \begin{discorso}
    \item Uso due reti di sei agenti: un ciclo diretto, in cui ogni agente ascolta solo il successivo, e un anello bidirezionale, in cui ogni agente ascolta i due vicini con peso un mezzo.
    \item {\color{blue}In entrambe cambio solo la riga dell'agente 1, con un parametro $t$ fra 0 e 1. L'agente 1 dà peso $t$ a se stesso, cioè alla propria opinione attuale, e moltiplica per $1-t$ i pesi che dava agli altri nella rete originale. Le altre righe non cambiano.}
    \item {\color{blue}Sul ciclo la riga dell'agente 1 diventa $t$ su se stesso e $1-t$ sull'agente 2. Sull'anello diventa $t$ su se stesso e $\frac{1-t}{2}$ su ciascuno dei due vicini, gli agenti 2 e 6.}
    \item {\color{blue}Nei grafi gli archi senza etichetta pesano 1 sul ciclo e un mezzo sull'anello. In particolare, sull'anello gli agenti 2 e 6 ascoltano l'agente 1 con peso un mezzo per ogni $t$: le loro righe non cambiano.}
    \item {\color{blue}Quindi $t$ è il self-weight dell'agente 1: con $t=0$ ho le reti originali, con $t=1$ l'agente 1 è stubborn.}
    \item {\color{blue}La condizione iniziale, salvo dove lo dico, è la stessa: opinioni da 0 a 1 con passo 0.2, dall'agente 1 all'agente 6.}
    \item {\color{blue}\clic Ora i casi di $t$. La tabella in basso vale per entrambe le reti.}
    \item Per $t=0$ entrambe sono fortemente connesse e periodiche: sul ciclo l'unico ciclo ha lunghezza 6, sull'anello ogni ciclo ha lunghezza pari, quindi periodo 6 e periodo 2. E sono doubly-stochastic: l'autovettore sinistro dell'autovalore 1 è uniforme, nessun agente è favorito.
    \item {\color{red}\barrato{Cambio} \barrato{solo} \barrato{la} \barrato{riga} \barrato{dell'agente} \barrato{1:}} {\color{red}$t$} {\color{red}\barrato{è} \barrato{il} \barrato{self-weight} \barrato{dell'agente} \barrato{1,} \barrato{il} \barrato{peso} \barrato{che} \barrato{dà} \barrato{alla} \barrato{propria} \barrato{opinione} \barrato{attuale.} \barrato{Si} \barrato{va} \barrato{dalle} \barrato{reti} \barrato{originali,}} {\color{red}$t=0$,} {\color{red}\barrato{all'agente} \barrato{1} \barrato{stubborn,}} {\color{red}$t=1$.}
    \item {\color{red}\barrato{La} \barrato{tabella} \barrato{in} \barrato{basso} \barrato{vale} \barrato{per} \barrato{entrambe} \barrato{le} \barrato{reti:} \barrato{per} \barrato{ogni}} {\color{red}$t$} {\color{red}\barrato{hanno} \barrato{le} \barrato{stesse} \barrato{proprietà,} \barrato{cambia} \barrato{solo} \barrato{il} \barrato{periodo} \barrato{per}} {\color{red}$t=0$.}
    \item Per $t$ fra 0 e 1 restano fortemente connesse, perché l'agente 1 ascolta ancora i suoi vicini. Il self-loop dell'agente 1 è un ciclo di lunghezza 1: nessun intero maggiore di 1 divide tutte le lunghezze dei cicli, quindi sono aperiodiche. E non sono più doubly-stochastic: la colonna dell'agente 1 somma a $1+t$.
    \item Per $t=1$ l'agente 1 ascolta solo se stesso: {\color{blue}l'unico arco che esce da lui è il self-loop, quindi} da lui non {\color{red}\barrato{esce}} {\color{blue}parte} nessun {\color{red}\barrato{arco,} \barrato{quindi}} {\color{blue}cammino verso gli altri, e} le reti non sono più fortemente connesse. La periodicità è definita solo per un grafo fortemente connesso, per questo nella tabella c'è un trattino: per $t=1$ conta il sink, lo vediamo nella prossima slide.
    \item {\color{red}\barrato{La} \barrato{condizione} \barrato{iniziale} \barrato{è} \barrato{sempre} \barrato{la} \barrato{stessa:} \barrato{opinioni} \barrato{da} \barrato{0} \barrato{a} \barrato{1,} \barrato{con} \barrato{passo} \barrato{0.2.}}
    \item Vediamo cosa succede alle opinioni.
  \end{discorso}
\end{slidedue}

\begin{slide}{m-fdg-conv}{One self-weight brings agreement}{1:10}
  \begin{discorso}
    \item Le colonne sono $t=0$, un mezzo e 1; sopra il ciclo, sotto l'anello.
    \item Per capire perché, uso il teorema del consensus. Ha tre condizioni equivalenti; qui la più comoda è quella sul grafo: esiste un nodo globally reachable, e il sottografo dei nodi globally reachable è aperiodico.
    \item Per $t=0$ ogni nodo è globally reachable, ma il sottografo è tutta la rete, che è periodica: la condizione non vale. È il caso della Rosetta stone di una rete fortemente connessa e periodica: le opinioni oscillano per sempre, nessuna convergenza. Sullo spettro ci sono autovalori diversi da 1 sul cerchio unitario.
    \item Basta un self-weight $t$ positivo, di un solo agente, e il gruppo raggiunge il consensus.
    \item Per $t$ fra 0 e 1 ogni nodo è ancora globally reachable, e la rete, con il self-loop dell'agente 1, è aperiodica: la condizione vale, consensus. È il caso della rete fortemente connessa e aperiodica: $A(t)$ è primitiva.
    \item Per $t=1$ l'agente 1 è stubborn ed è l'unico nodo globally reachable. Sul ciclo gli altri formano una catena che finisce in lui; sull'anello gli agenti da 2 a 6 si ascoltano tra loro, e il 2 e il 6 ascoltano l'agente 1. Il sottografo del solo agente 1 ha il self-loop, quindi è aperiodico: la condizione vale ancora.
    \item Detto con il condensation digraph: un solo sink, l'agente 1, aperiodico. È il caso della Rosetta stone con un unico sink aperiodico: tutti finiscono sulla sua opinione iniziale, zero.
    \item Lo spettro dice la stessa cosa: per ogni $t$ positivo gli autovalori diversi da 1 sono strettamente dentro il disco unitario.
    \item Questo è il primo punto del lavoro: la testardaggine ribalta la dinamica. Un agente che dà peso alla propria opinione attuale, senza essere stubborn, basta perché su una rete fortemente connessa tutto il gruppo si metta d'accordo.
    \item Con il consensus possiamo misurare il peso di ciascun agente.
  \end{discorso}
\end{slide}

\begin{slide}{m-fdg-sp}{Social power grows gradually}{0:45}
  \begin{discorso}
    \item Per $t$ in $(0,1]$ c'è consensus, e la social power ha la forma chiusa a destra, la stessa sulle due reti.
    \item In alto: la social power dell'agente 1 cresce gradualmente con il suo self-weight, da $1/6$, che è il limite per $t$ che tende a 0, fino a 1 per $t=1$. A $t=\tfrac12$ vale $2/7$, e ogni altro agente $1/7$.
    \item In basso il valore di consensus, $3(1-t)/(6-5t)$: da 0.5 nel limite per $t$ che tende a 0, fino a $x_1(0)=0$; a $t=\tfrac12$ vale $3/7$.
    \item I valori sono gli stessi sulle due reti: la rete decide come le opinioni convergono, non dove.
    \item In termini sociali: l'inerzia guadagna influenza poco alla volta, e per $t$ fra 0 e 1 ogni opinione iniziale entra ancora nel consensus.
    \item Perché la stessa social power su due reti così diverse?
  \end{discorso}
\end{slide}

\begin{slidedue}{m-fdg-why}{Why the same social power}{1:20}
  \begin{discorso}
    \item Perché la stessa social power su due reti diverse? Parto dalla regola che la definisce: l'equazione $w^\top A=w^\top$, scritta per l'agente $j$.
    \item Il peso dell'opinione iniziale di $j$ nel consensus è la somma, su tutti gli agenti che ascoltano $j$, del peso con cui lo ascoltano per il peso della loro opinione iniziale. Il motivo: dopo un passo, $x_j(0)$ si trova solo nelle opinioni di chi ascolta $j$, e il consensus raggiunto partendo dalle opinioni al passo 1 è lo stesso.
    \item A $t=0$, in entrambe le reti, ogni colonna di $A(0)$ somma a uno: per ogni agente, i pesi con cui gli altri lo ascoltano sommano a uno. Nel ciclo diretto ognuno è ascoltato solo dal precedente, con peso uno; nell'anello dai due vicini, con peso un mezzo ciascuno.
    \item Quindi pesi tutti uguali soddisfano la regola, in tutte e due le reti: la struttura non favorisce nessuno. Ma a $t=0$ non c'è consensus: il vettore uniforme non è una social power, è solo il punto di partenza, il limite di $w(t)$ per $t$ che tende a zero.
    \item Per $t$ positivo cambia solo la riga dell'agente 1, nello stesso modo nelle due reti: peso $t$ alla propria opinione, $1-t$ alle opinioni che ascoltava, nelle stesse proporzioni di prima.
    \item \clic Supponiamo che gli agenti da 2 a 6 abbiano tutti lo stesso peso $b$.
    \item L'agente 1 è ascoltato dagli stessi agenti di prima, con peso totale uno, e in più mantiene peso $t$ sulla propria opinione: $w_1=t\,w_1+b$. Quindi $w_1$ è maggiore di $b$, e il valore di consensus si sposta verso la sua opinione iniziale.
    \item Gli agenti che l'agente 1 non ascolta sono ascoltati come a $t=0$: ricevono $b$.
    \item Quelli che l'agente 1 ascolta ricevono da lui la quota $(1-t)\,a_{1j}\,w_1$, che è uguale ad $a_{1j}\,b$ perché $(1-t)\,w_1=b$: la stessa quota che a $t=0$ ricevevano da un agente di peso $b$. Ricevono $b$ anche loro: l'ipotesi è coerente.
    \item In tutto questo si usano solo due fatti: le colonne di $A(0)$ sommano a uno, e la riga dell'agente 1 cambia nello stesso modo. Quanti agenti ascoltano un agente, e in che verso, non entra: l'agente 3 riceve $b$ dal solo agente 2 nel ciclo, e metà $b$ dall'agente 2 e metà dall'agente 4 nell'anello.
    \item La rete decide come le opinioni convergono, non dove.
    \item In fondo la verifica del report: con la sostituzione tutte le equazioni si riducono a $b=(1-t)\,w_1$, e con la somma uguale a uno si ritrova $w(t)$, unica per $t$ positivo.
    \item Finora al massimo un agente stubborn. Vediamo cosa succede con due.
  \end{discorso}
\end{slidedue}

\begin{slide}{m-bridge}{Two stubborn agents}{1:00}
  \begin{discorso}
    \item Come nell'esempio con $t=1$, l'agente 1 è stubborn; ora sull'anello lo è anche l'agente 4, quello opposto.
    \item Il condensation digraph ha due sink, $\{1\}$ e $\{4\}$, aperiodici per i self-loop.
    \item Per il teorema del disagreement il gruppo converge a un disagreement: gli agenti 2, 3, 5 e 6 finiscono su combinazioni convesse delle opinioni iniziali di 1 e di 4.
    \item Le righe del limite di $A^k$ sono diverse, quindi la social power non è definita.
    \item Qui mi fermo con il modello FDG, per due motivi: è già visto nel corso, e la social power, che è ciò che mi interessa, richiede il consensus.
    \item Però lascio aperta una domanda. Gli agenti 1 e 4 non cambiano mai opinione: sono membri del gruppo, o sono opinioni fisse, esterne, che gli altri ascoltano?
    \item E in generale: un agente la cui opinione resta invariata indipendentemente dalle opinioni degli altri è un membro del gruppo, o un'opinione fissa esterna che gli altri ascoltano?
    \item Ci torno più avanti. Ora il secondo modello, il modello di Friedkin–Johnsen.
  \end{discorso}
\end{slide}

\parte{Parte 3 — Il modello FJ, teoria}{6:55 {\color{red}\barrato{5:45}}}
\begin{slide}{m-fj-model}{The FJ model}{1:00}
  \begin{discorso}
    \item Nel modello FJ ogni agente pesa con $\lambda_i$ la media delle opinioni che ascolta, e con il complemento $1-\lambda_i$ il suo prejudice $u_i$.
    \item $W$ è row-stochastic, la matrice delle influenze sociali: ha il ruolo che $A$ aveva nel FDG.
    \item $\lambda_i$ è la susceptibility dell'agente $i$ all'influenza sociale; $u$ è il vettore costante dei prejudice, opinioni fisse che gli agenti ascoltano; $1-\lambda_i$ è l'attachment dell'agente al proprio prejudice.
    \item A destra i casi estremi di $\lambda_i$. Con $\lambda_i=1$ l'agente non è attaccato al suo prejudice e segue la regola del FDG; se vale per tutti, $\Lambda=I$, il modello è proprio il FDG con $A=W$. Per questo da qui in poi almeno un agente ha $\lambda_i<1$.
    \item Un agente con $\lambda_i<1$ si dice prejudiced: il suo prejudice entra davvero nell'aggiornamento.
    \item All'estremo opposto, $\lambda_i=0$: l'agente non ascolta più nessuno, e dal primo passo in poi la sua opinione è $u_i$. È stubborn solo se $u_i$ è la sua opinione iniziale; altrimenti salta su $u_i$ al primo passo e resta lì.
    \item Nel FJ originale il prejudice è l'opinione iniziale, $u=x(0)$, e l'attachment è il self-weight $w_{ii}$. Il tutorial non lo impone, e io nemmeno; negli esempi userò $u=x(0)$.
    \item Rispetto al FDG cambia ciò a cui l'agente è legato: il self-weight lo lega alla sua opinione attuale, che continua a muoversi verso gli altri; il prejudice lo lega a un valore fisso, che non si muove mai.
    \item Ora le stesse tre domande per il modello FJ. La prima: quando converge?
  \end{discorso}
\end{slide}

\begin{slide}{m-fj-conv}{Convergence in the FJ model}{0:55 {\color{red}\barrato{0:50}}}
  \begin{discorso}
    \item Il modello FJ converge se il limite esiste per ogni condizione iniziale; è asintoticamente stabile se $\Lambda W$ è convergente, cioè ha raggio spettrale minore di uno.
    \item {\color{red}\barrato{A} \barrato{sinistra:} \barrato{un}} {\color{blue}Entrambe le proprietà dipendono da quali agenti i prejudice raggiungono. Un} agente è P-dependent se è prejudiced, o se è influenzato, direttamente o attraverso altri, da un agente prejudiced. Altrimenti è P-independent.
    \item {\color{red}\barrato{A} \barrato{destra} \barrato{la} \barrato{condizione} \barrato{di} \barrato{stabilità} \barrato{asintotica:} \barrato{il} \barrato{modello} \barrato{FJ} \barrato{è} \barrato{asintoticamente} \barrato{stabile} \barrato{se} \barrato{e} \barrato{solo} \barrato{se} \barrato{tutti} \barrato{gli} \barrato{agenti} \barrato{sono} \barrato{P-dependent.} \barrato{Succede} \barrato{per} \barrato{esempio} \barrato{su} \barrato{una} \barrato{rete} \barrato{fortemente} \barrato{connessa} \barrato{con} \barrato{almeno} \barrato{un} \barrato{agente} \barrato{prejudiced.}}
    \item {\color{red}\barrato{In} \barrato{basso} \barrato{l'intuizione.} \barrato{Lo} \barrato{scarto} \barrato{dall'equilibrio}} {\color{red}$Vu$} {\color{red}\barrato{evolve} \barrato{come}} {\color{red}$e(k+1)=\Lambda W\,e(k)$:} {\color{red}$u$} {\color{red}\barrato{non} \barrato{compare.} \barrato{Quindi}} {\color{red}$\Lambda W$} {\color{red}\barrato{decide} \barrato{se} \barrato{le} \barrato{opinioni} \barrato{convergono,}} {\color{red}$u$} {\color{red}\barrato{decide} \barrato{dove.}}
    \item Un agente {\color{blue}P-independent non è} prejudiced, {\color{red}\barrato{a} \barrato{ogni} \barrato{passo,} \barrato{trattiene} \barrato{solo} \barrato{una} \barrato{frazione}} {\color{red}$\lambda_i$} {\color{red}\barrato{degli} \barrato{scarti} \barrato{che} \barrato{ascolta:} \barrato{gli} \barrato{scarti} \barrato{si} \barrato{spengono} \barrato{solo}} {\color{blue}quindi $\lambda_i=1$, e non è influenzato da nessun agente P-dependent:} se {\color{red}\barrato{ogni} \barrato{agente} \barrato{è}} {\color{blue}lo fosse, attraverso di lui sarebbe} influenzato da un agente prejudiced.
    \item {\color{blue}Poiché $\Lambda$ è diversa dall'identità, almeno un agente è prejudiced, quindi almeno uno è P-dependent: $r$ è almeno 1.}
    \item {\color{blue}Numero prima i P-dependent, da 1 a $r$, poi i P-independent, e divido $x$, $u$, $W$ e $\Lambda$ in blocchi allo stesso modo. Il blocco $W^{21}$ è zero, perché i P-independent non ascoltano i P-dependent, e $\Lambda^{22}$ è l'identità, perché i P-independent non sono prejudiced.}
    \item {\color{blue}Il modello si divide in due: le opinioni dei P-dependent dipendono dai loro prejudice e dalle opinioni dei P-independent; quelle dei P-independent evolvono da sole.}
    \item {\color{blue}$W^{22}$ dice come i P-independent si ascoltano tra loro, ed è row-stochastic, perché $W$ lo è e $W^{21}$ è zero. Quindi la seconda equazione è un modello FDG fra i P-independent, che nessun prejudice raggiunge; per i P-dependent, le loro opinioni sono un ingresso in più.}
    \item {\color{blue}Su questa divisione si basa il teorema di convergenza.}
    \item {\color{red}\barrato{E} \barrato{non} \barrato{serve} \barrato{nessuna} \barrato{aperiodicità,} \barrato{a} \barrato{differenza} \barrato{del} \barrato{FDG.}}
    \item {\color{red}\barrato{E} \barrato{se} \barrato{qualche} \barrato{agente} \barrato{è} \barrato{P-independent?}}
  \end{discorso}
\end{slide}

\begin{slide}{m-fj-conv2}{Convergence with P-independent agents}{0:45 {\color{red}\barrato{0:50}}}
  \begin{discorso}
    \item {\color{red}\barrato{Gli} \barrato{agenti} \barrato{P-independent} \barrato{non} \barrato{ascoltano} \barrato{nessun} \barrato{agente} \barrato{P-dependent.} \barrato{Mettendo} \barrato{prima} \barrato{i} \barrato{P-dependent,} \barrato{la} \barrato{matrice}} {\color{red}$W$} {\color{red}\barrato{si} \barrato{divide} \barrato{in} \barrato{blocchi,} \barrato{e} \barrato{il} \barrato{blocco} \barrato{in} \barrato{basso} \barrato{a} \barrato{sinistra} \barrato{è} \barrato{zero.}}
    \item {\color{red}\barrato{Le} \barrato{opinioni} \barrato{dei} \barrato{P-dependent} \barrato{dipendono} \barrato{dai} \barrato{prejudice} \barrato{e} \barrato{dalle} \barrato{opinioni} \barrato{dei} \barrato{P-independent;} \barrato{quelle} \barrato{dei} \barrato{P-independent} \barrato{evolvono} \barrato{da} \barrato{sole.}}
    \item {\color{red}\barrato{Il} \barrato{blocco}} {\color{red}$W^{22}$} {\color{red}\barrato{dice} \barrato{come} \barrato{i} \barrato{P-independent} \barrato{si} \barrato{ascoltano} \barrato{tra} \barrato{loro:} \barrato{è} \barrato{un} \barrato{modello} \barrato{FDG} \barrato{a} \barrato{sé,} \barrato{che} \barrato{nessun} \barrato{prejudice} \barrato{raggiunge.}}
    \item Il teorema {\color{red}\barrato{generale} \barrato{di} \barrato{convergenza,} \barrato{a} \barrato{destra,}} dice {\color{red}\barrato{due}} {\color{blue}tre} cose. La {\color{blue}prima: la} parte {\color{blue}dei} P-dependent è sempre asintoticamente {\color{red}\barrato{stabile.} \barrato{E} \barrato{il} \barrato{modello} \barrato{FJ} \barrato{converge} \barrato{se} \barrato{e} \barrato{solo} \barrato{se} \barrato{tutti} \barrato{sono} \barrato{P-dependent,} \barrato{oppure} \barrato{converge} \barrato{il} \barrato{modello} \barrato{FDG} \barrato{dei} \barrato{P-independent.}} {\color{blue}stabile, cioè $\Lambda^{11}W^{11}$ ha raggio spettrale minore di uno.}
    \item {\color{red}\barrato{Quando} \barrato{converge,} \barrato{le} \barrato{opinioni} \barrato{finali}} {\color{blue}La seconda: il modello FJ converge se e solo se tutti gli agenti sono P-dependent, oppure converge il modello FDG} dei {\color{red}\barrato{P-dependent} \barrato{sono} \barrato{combinazioni} \barrato{convesse} \barrato{dei} \barrato{prejudice} \barrato{e} \barrato{delle} \barrato{opinioni} \barrato{finali} \barrato{dei} \barrato{P-independent.}} {\color{blue}P-independent, cioè $W^{22}$ è regular.}
    \item {\color{blue}La terza: quando converge, le opinioni finali dei P-dependent sono $V$ per il vettore dei loro prejudice e delle opinioni finali dei P-independent. $V$ è stocastica: ogni opinione finale è una combinazione convessa dei prejudice e delle opinioni finali dei P-independent.}
    \item Quindi l'aperiodicità torna a servire, ma solo per i P-independent: {\color{red}\barrato{lo}} {\color{blue}il loro modello FDG converge se e solo se tutti i sink del loro condensation digraph sono aperiodici, come dice la Rosetta stone. Lo} vedremo in un esempio.
    \item {\color{blue}Il caso più semplice è quello senza P-independent.}
    \item {\color{red}\barrato{Seconda} \barrato{domanda:} \barrato{dove} \barrato{convergono} \barrato{le} \barrato{opinioni?}}
  \end{discorso}
\end{slide}

\begin{slide}{m-fj-conv3}{All agents P-dependent\nuova}{0:50}
  \begin{discorso}
    \item {\color{blue}Se tutti gli agenti sono P-dependent, $r=n$: non ci sono P-independent, e la prima equazione è tutto il modello.}
    \item {\color{blue}Il corollario dice che il modello FJ è asintoticamente stabile se e solo se tutti gli agenti sono P-dependent, e allora $V=(I-\Lambda W)^{-1}(I-\Lambda)$. Succede per esempio su una rete fortemente connessa con almeno un agente prejudiced, come le nostre due reti.}
    \item {\color{blue}Con $r=n$ la formula del teorema dà, per ogni condizione iniziale, $x(\infty)=Vu$.}
    \item {\color{blue}In basso l'intuizione. $Vu$ è un equilibrio del modello. Lo scarto dall'equilibrio, $e(k)=x(k)-Vu$, evolve come $e(k+1)=\Lambda W\,e(k)$: $u$ non compare. Quindi $\Lambda W$ decide se le opinioni convergono, $u$ decide dove.}
    \item {\color{blue}A ogni passo l'agente $i$ tiene una frazione $\lambda_i$ della media pesata degli scarti degli agenti che ascolta: un agente non prejudiced fa solo la media, un agente prejudiced ne perde una frazione $1-\lambda_i$. Gli scarti escono dal gruppo solo attraverso gli agenti prejudiced, e si spengono solo se ogni agente è influenzato, direttamente o attraverso altri, da uno di loro.}
    \item {\color{blue}E non serve nessuna aperiodicità, a differenza del FDG.}
    \item {\color{blue}Seconda domanda: dove convergono le opinioni?}
  \end{discorso}
\end{slide}

\begin{slide}{m-fj-final}{Final opinions and the augmented graph}{1:00 {\color{red}\barrato{1:10}}}
  \begin{discorso}
    \item Con tutti gli agenti P-dependent le opinioni finali sono $x(\infty)=Vu$.
    \item Tre conseguenze. Le opinioni iniziali sono dimenticate: $x(0)$ non compare. Contano solo i prejudice degli agenti prejudiced: se $\lambda_j=1$, la colonna $j$ di $V$ è nulla.
    \item E, per un dato $u$, il gruppo raggiunge un consensus se e solo se tutti gli agenti prejudiced hanno lo stesso prejudice.
    \item Il tutorial dà anche una seconda lettura, il modello augmented: a ogni agente $i$ si affianca un agente virtuale stubborn $i'$, che vale sempre $u_i$.
    \item La matrice $\bar A$ è row-stochastic, quindi il modello augmented è un modello FDG su $2n$ agenti. Nella figura l'agente 1 è l'unico prejudiced: a sinistra il modello FJ, a destra il grafo augmented con l'agente virtuale $1'$.
    \item {\color{blue}I grafi sono semplificati, per motivi di spazio. Nel modello FJ ogni agente ha un prejudice $u_i$, ma solo l'agente 1 gli dà peso: per gli altri $1-\lambda_i=0$, quindi disegno solo $u_1$. Il grafo augmented ha $2n=12$ nodi, e ne disegno 7, i sei agenti reali e $1'$: gli agenti virtuali da $2'$ a $6'$ non li ascolta nessuno, perché l'agente $i$ ascolta $i'$ con peso $1-\lambda_i$, che per gli agenti da 2 a 6 è zero. È la stessa scelta della figura del report.}
    \item Ogni agente virtuale è {\color{red}\barrato{un} \barrato{sink:}} {\color{blue}stubborn: l'unico arco che esce} da lui {\color{red}\barrato{non} \barrato{esce} \barrato{nessun} \barrato{arco,} \barrato{a} \barrato{parte}} {\color{blue}è} il suo {\color{red}\barrato{self-loop,} \barrato{che} \barrato{lo} \barrato{rende} \barrato{aperiodico.}} {\color{blue}self-loop. Quindi è un sink del grafo augmented, aperiodico per il self-loop.}
    \item {\color{red}\barrato{Con} \barrato{tutti} \barrato{gli} \barrato{agenti} \barrato{P-dependent} \barrato{sono} \barrato{gli} \barrato{unici} \barrato{sink:} \barrato{ogni} \barrato{agente} \barrato{reale} \barrato{ha} \barrato{un} \barrato{cammino} \barrato{verso} \barrato{un} \barrato{agente} \barrato{prejudiced,} \barrato{e} \barrato{da} \barrato{lì} \barrato{verso} \barrato{il} \barrato{suo} \barrato{agente} \barrato{virtuale,} \barrato{quindi} \barrato{nessun} \barrato{gruppo} \barrato{di} \barrato{agenti} \barrato{reali} \barrato{è} \barrato{chiuso.}}
    \item {\color{red}\barrato{Il} \barrato{teorema} \barrato{del} \barrato{disagreement} \barrato{chiede} \barrato{solo} \barrato{questo:} \barrato{almeno} \barrato{due} \barrato{sink,} \barrato{tutti} \barrato{aperiodici.} \barrato{Non} \barrato{serve} \barrato{che} \barrato{la} \barrato{rete} \barrato{sia} \barrato{fortemente} \barrato{connessa,} \barrato{anzi} \barrato{con} \barrato{due} \barrato{sink} \barrato{non} \barrato{lo} \barrato{è} \barrato{mai;} \barrato{e} \barrato{ogni} \barrato{sink} \barrato{è} \barrato{per} \barrato{definizione} \barrato{una} \barrato{componente} \barrato{fortemente} \barrato{connessa,} \barrato{quindi} \barrato{ha} \barrato{senso} \barrato{dire} \barrato{che} \barrato{è} \barrato{aperiodico.}}
    \item {\color{red}\barrato{Il} \barrato{modello} \barrato{augmented} \barrato{converge:} \barrato{gli} \barrato{agenti} \barrato{virtuali} \barrato{restano} \barrato{sui} \barrato{prejudice,} \barrato{e} \barrato{gli} \barrato{agenti} \barrato{reali} \barrato{finiscono} \barrato{su} \barrato{combinazioni} \barrato{convesse} \barrato{dei} \barrato{prejudice,} \barrato{che} \barrato{è} \barrato{proprio}} {\color{red}$Vu$.} {\color{red}\barrato{Il} \barrato{disagreement} \barrato{è} \barrato{quello} \barrato{del} \barrato{modello} \barrato{augmented:} \barrato{per} \barrato{un} \barrato{dato}} {\color{red}$u$} {\color{red}\barrato{gli} \barrato{agenti} \barrato{reali} \barrato{possono} \barrato{anche} \barrato{essere} \barrato{d'accordo,} \barrato{se} \barrato{tutti} \barrato{i} \barrato{prejudice} \barrato{sono} \barrato{uguali,} \barrato{come} \barrato{diceva} \barrato{il} \barrato{terzo} \barrato{punto.}}
    \item In termini sociali: un prejudice si comporta come un agente in più che non cambia mai opinione, messo accanto all'agente prejudiced.
    \item Letto al contrario, questo risponderà alla domanda lasciata aperta. Prima, la terza domanda: quanto pesa ciascun agente?
  \end{discorso}
\end{slide}

\begin{slide}{m-fj-centrality}{Influence centrality}{0:40}
  \begin{discorso}
    \item La social power richiede il consensus, che nel modello FJ non è una proprietà del modello: dipende dai prejudice. E il modello augmented ha almeno $n$ sink, quindi per $n\ge2$ non va mai al consensus.
    \item Serve un'altra misura. Come nel tutorial, prendo il modello FJ asintoticamente stabile e $u=x(0)$, l'ipotesi tradizionale della teoria FJ: è questa che dà un ruolo alle opinioni iniziali, $x(\infty)=Vx(0)$.
    \item La influence centrality $c_i$ è il peso medio dell'opinione iniziale dell'agente $i$ nelle opinioni finali del gruppo: la media della colonna $i$ di $V$.
    \item $V$ è stocastica, quindi $c$ è non negativo e somma a uno, come la social power.
    \item Ora torno alla domanda lasciata aperta.
  \end{discorso}
\end{slide}

\begin{slidedue}{m-back-question}{Back to the question}{1:45 {\color{red}\barrato{1:15}}}
  \begin{discorso}
    \item Riprendo la domanda lasciata aperta con il modello FDG. Sull'anello gli agenti 1 e 4 sono stubborn, non cambiano mai opinione: sono membri del gruppo, o opinioni fisse esterne che gli altri ascoltano?
    \item Il modello augmented trasforma un prejudice in un agente stubborn virtuale. Letto al contrario: un agente stubborn è un prejudice per gli agenti che lo ascoltano.
    \item Lo vedo sull'anello: tolgo gli agenti 1 e 4. L'agente 2 aggiorna come $x_2(k+1)=\tfrac12x_3(k)+\tfrac12x_1(0)$: è il modello FJ con $\lambda_2=\tfrac12$ e prejudice $u_2=x_1(0)$.
    \item Allo stesso modo l'agente 6 ha prejudice $x_1(0)$, e gli agenti 3 e 5 hanno prejudice $x_4(0)$, tutti con $\lambda_i=\tfrac12$. Le traiettorie sono le stesse del FDG.
    \item \clic Quindi la risposta è: valgono entrambe le letture.
    \item Come membri del gruppo, 1 e 4 sono agenti FJ con $\lambda_i=0$ e prejudice uguale alla loro opinione iniziale, cioè stubborn, e gli altri seguono la regola del FDG.
    \item Come opinioni fisse esterne, sono il prejudice degli agenti che li ascoltano, come qui.
    \item Le due letture danno le stesse traiettorie. Il disagreement dei due agenti stubborn è il disagreement di agenti ancorati a prejudice diversi.
    \item {\color{blue}Questo disagreement lo dà il teorema del disagreement, che chiede solo questo: almeno due sink, tutti aperiodici. Non serve che la rete sia fortemente connessa, anzi con due sink non lo è mai; e ogni sink è per definizione una componente fortemente connessa, quindi ha senso dire che è aperiodico.}
    \item {\color{blue}Nel FDG i sink sono i due agenti stubborn, aperiodici per i self-loop. Nel modello augmented i sink sono gli agenti virtuali, aperiodici per i self-loop; con tutti gli agenti P-dependent sono gli unici sink: ogni agente reale ha un cammino verso un agente prejudiced, e da lì verso il suo agente virtuale, quindi nessun gruppo di agenti reali è chiuso.}
    \item {\color{blue}Per lo stesso teorema, nel modello augmented gli agenti reali finiscono su combinazioni convesse dei prejudice, che è proprio $Vu$. Il disagreement è quello dei $2n$ agenti del modello augmented, perché gli agenti virtuali restano sui loro prejudice; gli agenti reali possono anche essere d'accordo, se tutti gli agenti prejudiced hanno lo stesso prejudice.}
    \item Stesse traiettorie: ma è lo stesso sistema? Lo vedremo negli esempi, sulle stesse due reti del FDG.
  \end{discorso}
\end{slidedue}

\parte{Parte 4 — Il modello FJ, esempi}{9:00 {\color{red}\barrato{8:10}}}
\begin{slide}{m-fj-graphs}{Example networks: directed cycle}{0:45 {\color{red}\barrato{0:50}}}
  \begin{discorso}
    \item Per gli esempi uso le stesse reti del FDG. Qui il ciclo {\color{red}\barrato{diretto;}} {\color{blue}diretto,} l'anello {\color{red}\barrato{si} \barrato{tratta} \barrato{allo} \barrato{stesso} \barrato{modo.}} {\color{blue}nella prossima slide.} Nel modello FJ $W$ è la rete senza self-loop, $A(0)$, e il prejudice di ogni agente è la sua opinione iniziale.
    \item A sinistra un solo agente. Nel FDG, figura (a), l'agente 1 dà peso $t$ alla propria opinione attuale: è il self-loop. Nel FJ, figura (b), lo stesso peso $t$ va sul suo prejudice $u_1$: è la freccia che entra nell'agente 1, che non è un arco del grafo.
    \item A destra tutti gli agenti. Nel FDG, figura (c), ogni agente ha self-loop 0.2 e dà 0.8 al successivo; nel FJ, figura (d), ogni agente ha il suo prejudice con peso 0.2, e dà 0.8 al successivo.
    \item Stessi archi, stessi pesi: cambia solo l'oggetto del peso in più, l'opinione attuale nel FDG, un prejudice nel FJ.
    \item Non sto riscrivendo un modello nell'altro: costruisco la rete FJ analoga, per confrontare i due modelli.
    \item L'unica eccezione è $t=1$: la figura (b) diventa la riscrittura del tutorial, $\lambda_1=0$ e $u_1=x_1(0)$, con le stesse traiettorie della figura (a). È la risposta alla domanda.
    \item La rete (d) la ritroveremo più avanti, scritta come modello augmented.
    \item {\color{blue}Ora le stesse quattro reti sull'anello.}
    \item {\color{red}\barrato{Nella} \barrato{prossima} \barrato{slide} \barrato{le} \barrato{stesse} \barrato{cose} \barrato{in} \barrato{formule.}}
  \end{discorso}
\end{slide}

\begin{slide}{m-fj-graphs-ring}{Example networks: bidirectional ring\nuova}{0:35}
  \begin{discorso}
    \item {\color{blue}Le stesse quattro reti sull'anello bidirezionale, dove ogni agente ascolta i due vicini.}
    \item {\color{blue}A sinistra un solo agente. L'agente 1 dà peso $t$ alla sua opinione attuale nel FDG, figura (a), o al suo prejudice nel FJ, figura (b), e ascolta gli agenti 2 e 6 con peso $\frac{1-t}{2}$ ciascuno. Tutti gli altri archi pesano un mezzo: in particolare gli agenti 2 e 6 ascoltano l'agente 1 con peso un mezzo, per ogni $t$.}
    \item {\color{blue}A destra tutti gli agenti: ognuno dà peso 0.2 alla sua opinione attuale nel FDG, figura (c), o al suo prejudice nel FJ, figura (d), e 0.4 a ciascuno dei due vicini, cioè 0.8 diviso a metà.}
    \item {\color{blue}Valgono le stesse osservazioni del ciclo: stessi archi, stessi pesi, cambia solo l'oggetto del peso in più.}
    \item {\color{blue}Nella prossima slide le stesse cose in formule, per le due reti.}
  \end{discorso}
\end{slide}

\begin{slide}{m-fj-networks}{FDG and FJ on the same networks}{0:40 {\color{red}\barrato{0:35}}}
  \begin{discorso}
    \item La tabella dice in formule le figure di prima, per le due reti: $W$ è la rete originale, senza self-loop; stesse opinioni iniziali, e il prejudice di ogni agente è la sua opinione iniziale, come nel modello FJ originale.
    \item Confronto sempre i due modelli a parità di peso: nel FDG il peso va sull'opinione attuale, nel FJ sul prejudice.
    \item Prima riga, il caso estremo: l'agente 1 stubborn è $a_{11}=1$ nel FDG, $\lambda_1=0$ con $u_1=x_1(0)$ nel FJ.
    \item Seconda riga, {\color{red}\barrato{un} \barrato{solo} \barrato{agente}} {\color{blue}l'agente 1} con peso {\color{red}$t$.} {\color{blue}$t$, scritta per le due reti: sul ciclo l'agente 1 ascolta l'agente 2 con peso $1-t$, sull'anello la media degli agenti 2 e 6, con peso $1-t$ in tutto.} Le due equazioni sono uguali tranne un termine: nel FDG l'agente 1 dà peso $t$ alla sua opinione attuale $x_1(k)$, nel FJ al suo prejudice $u_1$. Quindi $1-\lambda_1=t$.
    \item Terza riga, tutti gli agenti con lo stesso peso 0.2: nel FDG $A=0.2I+0.8W$, nel FJ $\Lambda=0.8I$. {\color{blue}Ogni arco pesa 0.8 sul ciclo e 0.4 sull'anello.}
    \item Nelle prime due righe gli altri agenti restano come nella rete originale.
    \item Stesso peso, oggetto diverso: l'opinione attuale, che si muove ancora verso gli altri, o il prejudice, che non si muove mai.
    \item Comincio dal caso estremo, l'agente stubborn.
  \end{discorso}
\end{slide}

\begin{slidedue}{m-state-input}{Stubborn agent: state or input}{1:30 {\color{red}\barrato{1:15}}}
  \begin{discorso}
    \item Prima riga della tabella: l'agente 1 stubborn. Il tutorial riscrive così un modello FDG con agenti stubborn: $\lambda_i=0$ e $u_i=x_i(0)$ per gli stubborn, $\lambda_i=1$ per gli altri.
    \item A sinistra il FDG con $A(1)$, a destra il FJ con $\lambda_1=0$: le traiettorie sono le stesse, come per la domanda.
    \item Ma la stabilità no. Il modello FDG converge ma non è mai asintoticamente stabile, perché una matrice row-stochastic ha sempre l'autovalore 1; il modello FJ invece è asintoticamente stabile, perché tutti gli agenti sono P-dependent.
    \item \clic Per vederlo perturbo l'opinione iniziale dell'agente 1: $x_1(0)=0.5$, con $u$ invariato.
    \item Nel FDG l'agente 1 tiene 0.5 e il gruppo lo {\color{red}\barrato{segue:}} {\color{blue}segue. Il limite è $x_1(0)$ per tutti, perché la social power è $e_1$: passa da 0 a 0.5, e} la perturbazione non viene mai dimenticata. {\color{red}\barrato{Nel} \barrato{FJ} \barrato{l'agente} \barrato{1} \barrato{torna} \barrato{su}} {\color{red}$u_1=0$} {\color{red}\barrato{al} \barrato{primo} \barrato{passo,} \barrato{e} \barrato{il} \barrato{limite} \barrato{è} \barrato{lo} \barrato{stesso,}} {\color{red}$Vu$.}
    \item {\color{blue}Nel FJ l'agente 1 torna su $u_1=0$ al primo passo, e il gruppo va allo stesso limite di prima. È il significato di $x(\infty)=Vu$: il limite contiene solo i prejudice, non le opinioni iniziali, quindi perturbare $x(0)$ con $u$ fisso non lo cambia.}
    \item {\color{blue}In formule: la differenza fra le due traiettorie evolve come $\delta(k+1)=\Lambda W\,\delta(k)$, la stessa ricorsione degli scarti, e va a zero perché $\Lambda W$ è Schur stabile. Nel FDG la stessa differenza evolve con $A(1)$ e tende a $w^\top\delta(0)=0.5$ per tutti.}
    \item Sulle matrici: la riga dell'agente 1 è $e_1^\top$ in $A(1)$ ed è zero in $\Lambda W$, quindi l'autovalore 1 di $A(1)$ diventa 0 in $\Lambda W$.
    \item In termini sociali: nel FDG la stubbornness è una memoria nello stato, mai dimenticata; nel FJ è un ingresso esterno, e lo stato dimentica $x(0)$.
    \item Ora la seconda riga: un solo agente con peso un mezzo.
  \end{discorso}
\end{slidedue}

\begin{slide}{m-one-prej}{A single prejudice}{0:50}
  \begin{discorso}
    \item Seconda riga della tabella, con $t=\tfrac12$: a sinistra il FDG con $A(\tfrac12)$, peso sull'opinione attuale; a destra il FJ con $\lambda_1=\tfrac12$, peso sul prejudice.
    \item La rete è fortemente connessa, quindi tutti gli agenti sono P-dependent e il modello FJ è asintoticamente stabile, anche se $G$ è periodico e non ha self-loop.
    \item Nel FJ $V=\uno e_1^\top$: tutto il gruppo finisce sul prejudice dell'agente 1, zero, per ogni attachment positivo, per quanto piccolo. Nel FDG invece il consensus è a $3/7$.
    \item Sulla misura del potere: $c=e_1$ per ogni $\lambda_1<1$, l'agente 1 prende tutta la influence centrality; nel FDG, a $t=\tfrac12$, la sua social power è $2/7$.
    \item Gli altri due valori dell'attachment non hanno bisogno di una figura nuova. Con attachment zero nessun agente ha un prejudice, e il modello FJ è il modello FDG con $A(0)$: la colonna $t=0$ della figura del FDG, che non converge. Con attachment uno l'agente 1 è stubborn: è l'esempio della slide precedente.
    \item È di nuovo la testardaggine che ribalta la dinamica: un solo agente ancorato, che influenza tutti, decide dove finisce tutto il gruppo, per quanto debole sia il suo attachment.
    \item E se tutti gli agenti hanno un prejudice?
  \end{discorso}
\end{slide}

\begin{slide}{m-all-prej}{Every agent prejudiced}{0:45}
  \begin{discorso}
    \item Terza riga della tabella: tutti gli agenti hanno lo stesso peso 0.2, nel FDG come self-weight, nel FJ come attachment. I prejudice invece sono tutti diversi: ognuno ha come prejudice la propria opinione iniziale.
    \item Nel FDG $A$ è primitiva, perché la rete è fortemente connessa e ha self-loop, e doubly-stochastic: consensus sulla media, 0.5.
    \item Nel FJ tutti sono prejudiced, quindi il modello è asintoticamente stabile, e $V$ ha tutte le entrate positive: ogni prejudice entra in ogni opinione finale. I prejudice sono diversi: disagreement.
    \item Su una rete fortemente connessa il FDG non va mai in disagreement: nel FJ il disagreement viene dai prejudice.
    \item In termini sociali: l'inerzia alla fine cede, l'ancora no; ancore diverse tengono il gruppo diviso, quello che il tutorial chiama un cleavage di opinioni.
    \item La slide seguente mostra perché questo è un disagreement anche nel senso del corso.
  \end{discorso}
\end{slide}

\begin{slidetre}{m-six-sinks}{Six prejudices, six sinks}{1:30}
  \begin{discorso}
    \item È la rete (d), quella della Figura 10, scritta come modello FDG su dodici agenti: il modello augmented, con la stessa disposizione. Ogni prejudice $u_i$ diventa l'agente virtuale $i'$, stubborn; ogni agente reale dà 0.8 agli altri agenti reali e 0.2 al suo agente virtuale.
    \item Gli agenti reali formano una sola componente fortemente connessa, quella ombreggiata, che non è un sink. Collassandola si ottiene il condensation digraph: i sei agenti virtuali sono sei sink, aperiodici per i self-loop.
    \item Più sink, tutti aperiodici: è il caso del teorema del disagreement, come con più agenti stubborn, e il modello augmented converge. Lo stesso vale sull'anello.
    \item \clic Quindi $x(\infty)=Vu$, con $V$ a entrate positive: ogni opinione finale è una combinazione convessa di tutti e sei i prejudice.
    \item Anche con un solo agente prejudiced c'erano sei sink, ma solo $1'$ era ascoltato: togliendo gli altri cinque resta un solo sink, e il consensus.
    \item Così si chiude il filo del disagreement del FDG: un agente stubborn dà il consensus; due danno il disagreement, l'esempio della domanda; sei agenti virtuali, tutti ascoltati, danno il disagreement di questa slide.
    \item \clic Una nota sulle traiettorie. La figura mostra solo i sei agenti reali; nel modello augmented le traiettorie sarebbero dodici, e quelle dei sei agenti virtuali sono rette orizzontali sui prejudice.
    \item Ogni agente reale parte sulla retta del suo agente virtuale e la lascia appena la media delle opinioni che ascolta è diversa: sul ciclo subito, sull'anello gli agenti da 2 a 5 più tardi, quando arriva il cambiamento degli agenti 1 e 6.
    \item Il peso 0.2 sul suo agente virtuale lo richiama sempre, e si ferma su $Vu$, mai sul proprio prejudice. Nel FDG lo stesso 0.2 sta sull'opinione attuale, e l'agente segue il gruppo fino al consensus.
    \item Ultimo caso: dove nessun prejudice arriva.
  \end{discorso}
\end{slidetre}

\begin{slidedue}{m-aux}{Where no prejudice arrives}{1:10}
  \begin{discorso}
    \item Ultimo caso, escluso finora: non tutti gli agenti sono P-dependent.
    \item La rete: gli agenti 1, 2 e 3 formano un gruppo chiuso, un sink del condensation digraph; 4, 5 e 6 sono una catena che lo ascolta. In ogni esempio un solo agente è prejudiced, con $\lambda=0.5$; ombreggiati gli agenti P-dependent.
    \item Esempio 1: il prejudiced è l'agente 2, dentro il gruppo chiuso, che è un ciclo periodico. Tutti raggiungono l'agente 2, quindi tutti sono P-dependent. Ci aspettiamo stabilità asintotica, e tutti sul prejudice dell'agente 2, anche se il ciclo è periodico.
    \item Esempio 2: il prejudiced è l'agente 5, fuori dal gruppo chiuso. Il gruppo chiuso è P-independent: un modello FDG a sé, ed è periodico. Ci aspettiamo che il modello non converga.
    \item Esempio 3: come il secondo, più l'arco da 2 a 1, che rende il gruppo chiuso aperiodico. Ci aspettiamo convergenza, ma non stabilità asintotica: il limite deve dipendere dalle opinioni iniziali del gruppo chiuso. Per verificarlo, nel pannello (d) cambio solo $x_2(0)$.
    \item \clic Esempio 1, in alto a sinistra: tutti vanno a 0.5, il prejudice dell'agente 2.
    \item Esempio 2, in alto a destra: il gruppo chiuso oscilla, e il modello non converge.
    \item Esempio 3, in basso: converge; 1, 2, 3 e 6 vanno a 0.4, 4 e 5 a 0.65. Portando $x_2(0)$ da 0.5 a 1 diventano 0.6 e 0.75: il limite dipende davvero dalle opinioni iniziali dei P-independent.
    \item In termini sociali: gli agenti che nessun prejudice raggiunge formano un modello FDG a sé; hanno bisogno dell'aperiodicità e ricordano le loro opinioni iniziali.
    \item Infine, le due misure di potere.
  \end{discorso}
\end{slidedue}

\begin{slidedue}{m-lemma}{Centrality meets social power}{1:15}
  \begin{discorso}
    \item Ho due misure per la stessa domanda: quanto pesa l'opinione iniziale di ciascun agente.
    \item La social power, nel FDG: la quota di ciascun agente nell'opinione su cui il gruppo si mette d'accordo. Richiede il consensus.
    \item La influence centrality, nel FJ: quanto pesa in media l'opinione iniziale di ciascuno in opinioni finali che possono essere diverse. Richiede la stabilità asintotica e $u=x(0)$.
    \item Sono misure diverse: la social power pesa lo stato iniziale di un modello che ricorda $x(0)$, la influence centrality l'ingresso di un modello che lo dimentica. Per questo non le confronto direttamente, ma in parallelo.
    \item Si incontrano nel lemma a destra: se $W$ è fully regular, cioè il FDG con $A=W$ raggiunge il consensus, e tutti hanno la stessa susceptibility $\alpha$, per $\alpha$ che tende a 1 la influence centrality tende alla social power del FDG.
    \item Per vederlo prendo $W=A(\tfrac12)$, la rete con self-weight un mezzo sull'agente 1: è primitiva, quindi fully regular. E $u=x(0)$.
    \item Cosa ci aspettiamo. Per $\alpha=0$ ogni agente tiene la sua opinione iniziale: $V=I$, e la influence centrality è uniforme, $1/6$. Per $\alpha$ che tende a 1, $c$ tende a $w$: $2/7$ per l'agente 1, $1/7$ per gli altri.
    \item \clic Ed è quello che succede: $c_1$ sale da $1/6$ fino a $2/7$, $c_2$ scende fino a $1/7$. Sul ciclo $c_2$ prima scende anche sotto $1/7$.
    \item In termini sociali: man mano che gli agenti si aprono agli altri, l'inerzia dell'agente 1 conta di più.
    \item La social power dice chi guida un gruppo che si mette d'accordo; la influence centrality dice chi pesa di più in un gruppo che resta in disagreement. La seconda diventa la prima quando le ancore si allentano.
    \item Per concludere, i due modelli a confronto.
  \end{discorso}
\end{slidedue}

\parte{Parte 5 — Conclusioni}{1:30}
\begin{slide}{m-table}{FDG and FJ side by side}{0:40}
  \begin{discorso}
    \item La tabella mette a confronto i due modelli riga per riga.
    \item La resistenza: nel FDG è un self-weight, un'inerzia; nel FJ un attachment, un'ancora.
    \item L'agente stubborn: nel FDG è una condizione sulla sola matrice, $a_{ii}=1$; nel FJ sulla matrice e sull'ingresso, $\lambda_i=0$ e $u_i=x_i(0)$. La sua opinione è uno stato nel FDG, un ingresso nel FJ.
    \item La stabilità asintotica: mai nel FDG; nel FJ se e solo se tutti gli agenti sono P-dependent.
    \item Su una rete fortemente connessa il FDG dà consensus o nessuna convergenza, mai disagreement; il FJ dà disagreement ogni volta che i prejudice che contano sono diversi.
    \item L'aperiodicità serve sempre nel FDG; nel FJ solo per gli agenti P-independent.
    \item E la misura del potere: la social power richiede il consensus; la influence centrality richiede la stabilità asintotica e $u=x(0)$.
  \end{discorso}
\end{slide}

\begin{slide}{m-messages}{Take-home messages}{0:50}
  \begin{discorso}
    \item Tre messaggi da portare a casa.
    \item Primo, inerzia e ancora. Un self-weight rallenta un agente ma lo lascia seguire gli altri: su una rete fortemente connessa e periodica basta un self-weight fra 0 e 1 per trasformare le oscillazioni in consensus. Un prejudice non cede mai: un solo agente prejudiced che influenza tutti decide dove finisce il gruppo, e prejudice diversi lo tengono diviso.
    \item Secondo, stato e ingresso. La stessa opinione stubborn è uno stato nel modello FDG, che non è mai asintoticamente stabile, e un ingresso nel modello FJ, asintoticamente stabile quando tutti gli agenti sono P-dependent. Il grafo augmented riporta gli ingressi a stati.
    \item Terzo, il potere. La social power cresce gradualmente con il self-weight; un agente prejudiced che influenza tutti prende tutta la influence centrality; e con la stessa susceptibility per tutti e $W$ fully regular, la influence centrality tende alla social power quando $\alpha$ tende a 1.
  \end{discorso}
\end{slide}

\begin{slidepag}{\the\numexpr\pg{m-messages}+1\relax}{Thank you!}{}
  \begin{discorso}
    \item Grazie per l'attenzione. Sono a disposizione per le domande.
  \end{discorso}
\end{slidepag}

\end{document}
